\chapter{CONCLUSIONS}\label{ch:conclusion}
\section{Conclusions}
This laboratory exercise successfully demonstrated the practical application of Kanban principles in team task management. Through the paper-based system, team members experienced firsthand the benefits of workflow visualization, parallel execution, and peer verification. All nine tasks were completed with documented evidence across five board states.

The activity revealed both strengths and weaknesses in team coordination. While the visual system enhanced transparency and accountability, challenges emerged in task clarity and data consistency. The metaphorical "Mount Everest" task particularly emphasized the importance of clear requirements in collaborative work.

Key takeaways include the value of physical visualization in team settings, the efficiency gains from parallel task execution, and the critical role of verification in maintaining quality. These insights are directly transferable to software development contexts where agile methodologies and visual management tools are increasingly prevalent.

Future implementations would benefit from clearer initial task definitions, standardized recording methods, and more structured verification protocols. Overall, this exercise provided valuable experiential learning in collaborative workflow management and team coordination strategies.

%\section{Future Prospects of Our Work}
